\chapter{Future Work}
\label{chap:future-work}

The evaluation in \Cref{chap:validation} answers the research questions of this thesis, but
also identifies several directions for further work. These concern the use of the
knowledge-graph enrichment, wider repair coverage, comparison with a repair
baseline, and a broader human evaluation.

\section{Building on the Knowledge-Graph Enrichment}
\label{sec:fw-kg}

Objectives O4 and O6 are both met for the final evaluation, by two
distinct artefacts that describe the same repair attempts. The
\texttt{repair\_attempts} table of the final Gemma run holds one row per
repair-agent invocation, 235 in all (187 Round-1 and 48 Round-2 rows),
with the full explanation and PyPI evidence of every invocation. That
table, with its CSV export, is the structured benchmark dataset of O4.
The existing RML mapping, run over that export, materialised the same 235
attempts as 3{,}478 triples, each linked to the notebook resource of the
FAIR Jupyter knowledge graph that identifies the repaired notebook, and
the result passed every validation check in union with the existing
notebook graph (\Cref{sec:v-i6-fullscale}). That repair graph is the
enrichment of O6 at evaluation scale; it has not yet been loaded into the
public FAIR Jupyter triple store.

The first extension is deployment. The enriched triples exist as files
next to the run that produced them. Loading them into the FAIR Jupyter
triple store, beside the graph the upstream pipeline already publishes,
would make repair history queryable at the same public SPARQL endpoint as
the execution records it extends. Because the repair mapping is one more
\texttt{.rml.ttl} file of the shape the upstream build script already
iterates over, this requires running that script with the mapping and
the export in place, not new tooling. Materialising a second run next to
the first, such as the Qwen run of \Cref{sec:v-model-sensitivity}, would
first require the repair-attempt IRI template to carry the run identifier
as well as the row id, because each run's table numbers its rows from
one.

The second is richer use of the graph. The checks of
\Cref{sec:v-i6-fullscale} ask structural questions: does every attempt
link to a notebook, and does every executed attempt carry an outcome. The
same graph can also answer the questions for which the summary tables of
\Cref{chap:validation} were computed separately: which notebooks received
a Round-2 attempt and after which Round-1 outcome, which failing modules
abstained, and how the attempts of one run compare with those of another
once both are loaded. This is possible because the round, the model, and
the run identifier are recorded on every attempt. Written as SPARQL
queries over the union of the upstream graph and the repair triples,
these questions could be re-asked without returning to the SQLite table,
and the repair history could be combined with the execution, module, and
repository facts the graph already holds about the same notebooks.

The third is a richer vocabulary. The mapping records the proposed
package, version, and PyPI evidence as literals on the repair attempt.
Modelling the installed distribution and its release as resources of
their own, and linking each repair attempt to them, would connect repair
provenance to package and version entities that other datasets could
share. A query could then ask, for instance, which notebooks were
repaired by pinning the same NumPy release. Repair actions this thesis
deferred, such as replacing an import (\Cref{sec:fw-coverage}), and any
further repair rounds would need no change to the mapping: its existing
predicates already carry the action value and the round number.

\section{Widening Repair Coverage}
\label{sec:fw-coverage}

The single largest limit on this system's repair coverage is the static
import-to-distribution mapping. It left 162 of 235 repair-agent
invocations abstaining before the language model was reached
(\Cref{sec:v-final-abstention}), and the model-sensitivity experiment
showed this limit is shared across both models tested
(\Cref{sec:v-model-sensitivity}). Extending the mapping table is the
obvious first step, but it does not remove the underlying problem: any
hand-curated table needs continued maintenance and cannot resolve a name
it has never been told about. A more durable approach would resolve import
names against a maintained external source, and fall back to abstention
only when that source genuinely cannot answer.

Two further limits are narrower but real. The compatibility-evidence
registry that supports \texttt{wrong\_version} repairs holds three
hand-checked entries (\Cref{tab:m-compat}), which is enough for this
dataset and no more. Deriving such evidence automatically from release
notes or from a project's own API history would generalise it. The
repair scope is pip-only by design (\Cref{sec:pd-scope}), which leaves
the 10 \texttt{system\_library} records unrepairable and rules out fixing
an API rename inside a notebook's own source. A related extension
concerns the repository's declared dependency files, which the pipeline
currently uses only to rebuild the baseline environment
(\Cref{sec:m-repair-nodeps}). They could additionally be shown to the
repair model as non-authoritative context, for instance to flag that a
proposed pin conflicts with a version the repository declares, while PyPI
evidence remains the only source from which a version may be drawn.

A third question the evaluation raised but could not answer is how much a
larger repair budget would help. \Cref{sec:v-case-study} shows a notebook
that exposed a third dependency error after two successful targeted
repairs. Whether a third or fourth round would have recovered it, or
whether notebooks of this kind have a long tail of dependency problems,
is unmeasured: no run beyond two rounds was attempted.

Two further points concern how the final runs were executed rather than
what they repaired. First, the commit-pinning mechanism of FixApplicator
was not in effect in the final evaluation environment
(\Cref{sec:v-final-manifest}), so the evaluated repositories were cloned
at their default-branch state. A rerun with the metadata database
reachable from the execution environment would test each repository at the
commit that originally failed, and would show whether any of the 187
outcomes depends on changes the repository's authors made since. Second,
the Gemma explainer lost 27 of 237 explanation records to timeouts of the
local, CPU-only model service (\Cref{sec:v-final-explanation}). A larger
timeout, or serving the same model on faster hardware, would separate how
much of that gap is the machine and how much is the model; the present
results cannot.

\section{Baseline Comparison}
\label{sec:fw-baseline}

This thesis reports how its own repair layer performs, not how it performs
relative to an alternative. The most informative comparison would be
against the same agent with PyPI grounding removed, so that the
contribution of grounding could be isolated from the contribution of the
model. Such a comparison could reuse the existing evaluation design
without changing the dataset or the metric definitions of
\Cref{chap:validation}. The ungrounded baseline has not been built and
the comparison has not been run, so this thesis makes no claim about
grounding's measured benefit.

\section{Extending the Human Evaluation}
\label{sec:fw-human}

The human evaluation reported in \Cref{sec:v-human-eval-results} is a
convenience sample of 15 participants, each rating six explanations from
a frozen pool of 12, all of them Gemma explanations of original failures.
It supports a descriptive picture only
(\Cref{sec:v-human-eval-limitations}). The Round-2 explanations of
newly exposed errors and the Qwen explanations were evaluated
automatically for completion and schema validity only, and rating them
with the same instrument is the most direct extension. Three further
extensions would each answer a question the present study cannot.

With an appropriate sampling and analysis design, a larger sample could
support inferential comparisons rather than descriptive claims. A
replication with expert raters, judging explanations
against each notebook's true root cause, would measure objective
correctness, which this study does not: Q6 captures only whether an
explanation struck a reader as correct. Showing participants an
alternative explanation, whether from a different model or from a non-LLM
template, would place these results on a comparative scale.

The free-text feedback points to one concrete improvement worth testing
directly. Several participants asked for more about the likely cause and
less restatement of the error message (\Cref{sec:v-human-eval-results}).
Whether a revised prompt raises the satisfaction item without harming the
others could be tested with the existing instrument and example pool.
